\documentclass[12pt,a4paper]{article}
\usepackage[utf8]{inputenc}
\usepackage[T1]{fontenc}
\usepackage[brazil]{babel}
\usepackage{lmodern}
\usepackage[margin=3cm]{geometry}
\usepackage{setspace}
\usepackage{booktabs}
\usepackage{longtable}
\usepackage{array}
\usepackage{enumitem}
\usepackage{csquotes}
\usepackage{fancyhdr}
\usepackage[hidelinks]{hyperref}
\onehalfspacing
\pagestyle{fancy}
\fancyhf{}
\fancyhead[L]{\footnotesize Pesquisa de satisfação --- PetCuida \textit{Alpha 3i}}
\fancyhead[R]{\footnotesize \thepage}
\newenvironment{turnos}{\begin{list}{}{\setlength{\leftmargin}{2.2cm}\setlength{\labelwidth}{2.0cm}\setlength{\labelsep}{0.2cm}\setlength{\itemsep}{0.45em}\setlength{\parsep}{0pt}}}{\end{list}}
\newcommand{\fala}[1]{\item[\textbf{#1}\hfill]}
\title{\textbf{Pesquisa de Satisfação e Percepção do Aplicativo PetCuida\\ versão \textit{Alpha 3i}}\\[0.4cm]\large Entrevista semiestruturada com Atanizia e participação de Galdina}
\author{Entrevistador: Gabriel\\Participantes: Atanizia e Galdina}
\date{Registro com data do arquivo: 17 de setembro de 2026}
\begin{document}
\maketitle
\thispagestyle{empty}
\begin{abstract}
Este documento formaliza o registro de uma entrevista semiestruturada realizada por Gabriel com Atanizia, mãe do entrevistador, e com intervenções atribuídas provisoriamente a Galdina. A sessão teve como objetivos compreender dificuldades relacionadas ao cuidado de animais, apresentar o protótipo PetCuida \textit{Alpha 3i} e coletar percepção sobre sua utilidade e facilidade de uso. A análise indica que a entrevistada reconhece a relevância da proposta, mas enfatiza que a plataforma não substitui a responsabilidade do tutor. O principal valor percebido está no acesso a transporte, atendimento veterinário, castração, triagem e apoio comunitário.\par
\vspace{0.3cm}\noindent\textbf{Palavras-chave:} satisfação; usabilidade percebida; bem-estar animal; entrevista semiestruturada; PetCuida; Alpha 3i.
\end{abstract}
\newpage
\tableofcontents
\newpage
\section{Identificação da coleta}
\begin{longtable}{p{5cm}p{8.5cm}}\toprule\textbf{Item} & \textbf{Descrição}\\\midrule
Objeto avaliado & Aplicativo PetCuida, versão \textit{Alpha 3i}.\\
Instrumento & Entrevista semiestruturada com apresentação guiada do protótipo.\\
Entrevistador & Gabriel.\\
Participantes & Atanizia; Galdina com participação a confirmar no áudio.\\
Áudio & Arquivo digital em formato \texttt{opus}, com aproximadamente 25 minutos e 31 segundos.\\
Transcrição & Versão semiliteral revisada, acompanhada de arquivo \texttt{srt}.\\
Análise & Análise qualitativa exploratória por categorias temáticas.\\
\bottomrule\end{longtable}
\section{Objetivos}
O objetivo geral foi verificar a percepção de potenciais usuárias sobre o problema de acesso a cuidados para animais e sobre a proposta funcional do PetCuida. Os objetivos específicos foram: identificar barreiras econômicas e logísticas; compreender a relevância de castração, transporte e atendimento; apresentar o fluxo do protótipo; registrar percepção de utilidade e facilidade; e identificar limites éticos e metodológicos da sessão.
\section{Procedimentos e tratamento dos dados}
A sessão combinou contextualização do problema, conversa sobre experiências locais, apresentação guiada do protótipo e coleta de feedback. Não se tratou de teste de usabilidade autônomo: a maior parte da navegação foi conduzida pelo entrevistador. A transcrição foi editada em nível de superfície, com normalização de pontuação, remoção de repetições sem valor proposicional e correção de erros evidentes de reconhecimento. Termos incertos foram mantidos entre colchetes. A atribuição automática dos falantes foi preservada como hipótese de trabalho e não deve ser tratada como confirmação definitiva.
\section{Transcrição formalizada}
\noindent\textit{Convenções:} E designa Gabriel; A designa Atanizia; G designa Galdina quando a atribuição é provisória; [termo incerto] indica passagem que requer conferência auditiva; intervenções sem conteúdo proposicional foram reduzidas.

- \textit{entrevista teve uma prévia antes da gravação.}
\begin{turnos}
\fala{E:} Pode continuar a conversa. 
\fala{A:} Continuar? 
\fala{E:} É.
\fala{A:} Não, mas aí o que ele quer te entrevistar é sobre isso. O que você acha, porque já tem o hospital, só que é em Belo Horizonte. Está entendendo?
\fala{G:} De ter a criação de animais, sem ter condições de cuidar. 
\fala{A:} É.
\fala{E:} Mas existe também problema crônico, igual a minha mãe está comentando sobre o caso da gata da Lara(vizinha do entrevistador), que do nada deu um problema que fez o olho da gata saltar para fora porque estava inchado. E aí quando foi fazer...
\fala{A:} Levou no veterinário.
\fala{E:} Levou no veterinário, foi R\$250,00 só para fazer a extração.
\fala{A:} Só para tirar o olho da gatinha.
\fala{E:} Entendeu.
\fala{G:} Entendi.
\fala{G:} Eu, particularmente, eu não tenho, nunca tive. O Gabriel(filho da entrevistada Galdina) teve, dois cachorros. Mas eu nunca tive porque, exatamente por isso.
\fala{G:} Porque eu acho que a criação de animais é igual o filho. Quando você pega, você tem que cuidar. Se você não tem como cuidar, então é melhor você não ter.
\fala{G:} O filho, se você teve, você vai cuidar, indiferente de qualquer coisa. Mas foi uma escolha. E a criação de animais é da mesma forma.
\fala{G:} É uma escolha você ter. Se você precisa do médico, então tem que levar. Você tem que estar disposto a gastar.

- \textit{Observação: a senhora Galdina é baixa renda e trabalha como enfermeira na rede público-particular em Belo Horizonte.}

\fala{G:} Trabalha uma médica comigo que uma injeção é R\$16.000,00. R\$16.000,00. O que que ela fez?
\fala{G:} Ela pegou empréstimo, ela trabalha dobrado, que ela tem que fazer essa injeção contínua. Uma vez por ano ela tem que administrar essa injeção no cachorro. E ela faz.
\fala{G:} Ela gasta mais de R\$1.500,00 com cachorro. É convênio. Tudo que ela paga.
\fala{G:} Ração cara, ela compra carne para dar para os bichos. Aí, ela tem condições. Mas hoje, como que eu vou pegar um cachorro que eu não tenho condições de cuidar dele?
\fala{A:} Mas é o que a gente mais vê aqui no bairro. É isso. As pessoas não têm condições e deixam o cachorro ir criando a revelia.
\fala{G:} É.
\fala{A:} E depois o animal fica doente e elas não tem condições de criar. Que não é o nosso caso aqui. Essa gata estava abandonada lá na ONG.
\fala{E:} É.
\fala{A:} O Gabriel(entrevistador e filho da sra. Atanizia) ficou com dó e trouxe aqui para casa. Mas para que os outros não matassem. Ela(a gata resgatada) no caso, chegou até machucada.
\fala{E:} Ela tenho um defeito na boquinha. A minha boca é meio torta. Mas é por causa que ela machucou a mandíbula quando era pequena(1 mês de vida).
\fala{E:} Às vezes, quando ela vai ficar com a patinha assim(ao se assentar uma das patas não toca o chão). Mas isto veio de briga com outros gatos com ela(muitos gatos abandonados e não castrados).
\fala{G:} E no caso, esse trabalho para procurar uma solução para resolver esse problema? 
\fala{E:} Propôr. 
\fala{G:} O que tinha que propor é o quê?
\fala{G:} Castrar mesmo os animais que estão na rua. Para que não fique criando mais... 
\fala{A:} A castração para o pessoal daqui do bairro é difícil.
\fala{E:} É caro porque...
\fala{G:} Não, mas teria que ser tipo assim. Igual não tem o SAMU-PET de BH. O cachorro, o animal passou mal.
\fala{G:} Eles vão e buscam, similar ao que acontece em Belo Horizonte. Levantar não só para Belo Horizonte. Mas para todos os lugares.
\fala{G:} Porque o hospital para atendimento humano, tem só em Belo Horizonte? Não, tem nas cidades ao redor. Então, para as criações também terem ao redor.
\fala{G:} Porém, ter isso. A pessoa que não tem condições... É igual o filho.
\fala{G:} Não tem condições? Não tem. Criação não tem condições?
\fala{G:} Não tem. Porque se ter, cuidar. Agora, a pessoa ter consciência também disso.
\fala{G:} Porque eu ter um animal com a certeza de que é você que vai cuidar dele. É fácil. É a mesma coisa eu ter um filho.
\fala{G:} Com a certeza que você vai cuidar. A avó tem um neto. Vai cuidar.
\fala{G:} O filho vai fazer o quê? Vai arrumar outro. Porque quem está cuidando é a avó.
\fala{G:} É ter aquela responsabilidade com a criação de animais. Porque a criação de animais é a mesma coisa que um filho. É família.
\fala{G:} Entendeu? Então, assim... É conscientizar as pessoas de ter se for cuidar.
\fala{G:} Se não for cuidar, não ter. E, nesses casos, a pessoa está na rua... Ter uma ONG para estar realmente ajudando.
\fala{G:} Tipo nessa questão. Castrar, vacinar e afins...
\fala{A:} Acerca disto...
\fala{A:} Eu acho que a veterinária veio para atender a Tieta(cachorra idosa de 10 anos de uma vizinha).
\fala{A:} Por causa da patinha. Você viu ali que ela está até com o cone. A veterinária está vindo.
\fala{A:} Ela veio hoje.
\fala{E:} Ah, tá bom.
\fala{A:} Deve ser da prefeitura.
\fala{E:} Possivelmente.
\fala{A:} Porque a dona da Madalena(dona da Tieta), ela não tem condição.
\fala{E:} Não.
\fala{A:} Não pode pagar.
\fala{E:} Entendeu? Porque até mesmo, talvez, a Gabriela(filha da sra. Madalena)... 
\fala{A:} Sim.
\fala{E:} Ela até olha algumas coisas, mas não dá para ela ter.
\fala{A:} Não tem condições.
\fala{E:} Entendeu?
\fala{A:} É. Mas é isso. Entendeu?
\fala{G:} Então, assim... É ter mesmo para conscientizar as pessoas de que a criação de animais tem que ser igual ao filho. Até mesmo porque um cachorro, ele vai dando, dando cria, vai prejudicar o quê?
\fala{G:} A saúde da cachorra também. Então, é... Castrar.
\fala{G:} Aí... As pessoas têm que ter consciência.
\fala{E:} O problema que geralmente acontece, muitas vezes, que acontece aqui no contexto do bairro, é que geralmente o pessoal, assim, na questão de ter um afeto, às vezes, eles providenciam um animal. Até certo ponto da vida, o animal não dá problema nenhum. É igual o caso de Tieta.
\fala{E:} A Tieta deve estar com uns 10 anos já de vida, mais ou menos. Priscila(outra cachorra da sra. Madalena) viveu, basicamente, uns 12 anos sem nenhum problema. Só que, no fim da vida, vem as coisas que a pessoa não espera.
\fala{E:} É igual o caso dessa gatinha. A gente pegou ela, cuidamos o melhor que a gente pôde, só que a gente não esperava que, por exemplo, agora ela com, vamos dizer, 6 para 7 anos, né, porque ela nasceu em 2019, é... que ela ia ter esse problema por conta da idade e também por conta de briga com o gato, porque ela não aceita ser dominada.
\fala{E:} É igual, por exemplo, eu sou o dono dela, mas se ela não quer ficar em um canto, ela não fica. E ela tem raiva dos gatos, porque assim, ela briga, ela apanha, igual ela tem um problema aqui na patinha, igual você vai ver que ela fica assim, apoiada, e aí a patinha que ela tem machucada. Por causa que ela brigou com o gato.
\fala{E:} Porque o gato não respeita, e aí entra esse problema, o pessoal não consegue castrar o gato.
\fala{A:} Deixa o gato solto?
\fala{E:} Não, deixa o gato solto e o pessoal prefere mais arrumar gato(animal macho). Então, ou seja, teve uma época que ali para o lado da casa da dona Antônia(outra vizinha), tinha mais ou menos, eu acho que mais uns 20 gatos, né, e o problema que assim, muitos animais machos soltos poucas fêmeas.
\fala{A:} O volume de gatos em massa, sem castrar, e pouca fêmea, também sem castrar.
\fala{E:} E aí acaba que tem, igual aqui o caso, tem esse território todo aqui, né, que são basicamente dois, três quarteirão(150-300m por 250-300m). Mais ou menos uma média de cinco gatos a dez. E uma fêmea só, né, aí dá problema.
\fala{E:} E o problema é que eles ignoram o fato, por exemplo, que ela é castrada. Entendeu? Para a fêmea, né, já vai em cima.
\fala{E:} Aí é onde ela briga. Que é justamente isso. Essa pessoa, às vezes, até tem cuidado, tem condições de manter o animal, e vai cuidando mas...
\fala{E:} Só que acontece imprevistos, igual a esse caso, que ela brigou, possivelmente, deve ter, talvez, machucado, quebrado uma costela. Possivelmente, eu acho que deve ter, talvez, dado, trincado o osso, porque ela está com dificuldade de respirar. Só que não dá pra você levar ela num médico, porque justamente o custo é caro.
\fala{E:} Possivelmente, só a consulta dela, que era, talvez, uns 300, 500 reais. E aí, mais o tratamento. Então, possivelmente, entraria mais uns mil reais.
\fala{G:} Até o meio de transporte, né, que você não pode transportar de qualquer forma. Você tem que ter uma gaiola certa pra você levar. 
\fala{A:} Na gaiola. 
\fala{A:} Outra coisa, ali na casa do seu Francisco(outro vizinho), o pessoal estava abandonando as gatas, e os gatos iam pra lá, então as fêmeas iam tendo muitos filhotes.
\fala{A:} O que ele fez, juntamente com o filho Gilberto, para interromper a reprodução? Eles colocaram os gatos no carro e levaram em Belo Horizonte pra castrar, no hospital lá em Belo Horizonte, pra castrar. Mas nem todo mundo tem essa condição, não é?
\fala{A:} Então, de pegar o animal e levar. E os gatinhos, eles fizeram publicação na internet e doou os gatinhos. Ele falou que tem uma vez que ele estava com sete gatos na casa dele.
\fala{G:} É uma coisa, assim, que eu não gostaria de ter mesmo gato, porque eles não são animais que aceitam ser domesticados, que ficam ali dentro de casa mesmo. Eles ficam pra todos os lados. De qualquer lugar eles passam, eles sobem.
\fala{G:} Então, assim, é difícil. É difícil você conseguir chegar e ele tá lá dentro. Então, o cachorro ainda fica ali no terreiro só.
\fala{E:} O gato se relaciona com o problema que eu comentei sobre o caso da gatinha ali. Que é o seguinte, ele se permite ser domesticado, mas ele não é domesticado. Ou seja, ele vai pelo que é interessante pra ele.
\fala{G:} Sim. 
\fala{E:} E no caso do gato macho, tem vezes igual, por exemplo, o gato é daqui, mas a gente vai achar ele, por exemplo, 500 metros pra frente. Igual no caso ali, tinha gato que era daqui, da rua que a gente mora, no ponto de ônibus lá embaixo.
\fala{E:} Porque eles andam todos os dias. É. Né?
\fala{A:} É.
\fala{E:} E assim, em busca de alimento mesmo. 
\fala{A:} E de fêmea. 
\fala{E:} E de fêmea, principalmente.
\fala{E:} Já a gata fêmea é até melhor que o gato. Porque o que acontece? Ela é mais dócil.
\fala{G:} É?
\fala{E:} E também ela não tem a preocupação de sair do território dela. Então, por exemplo, ela(o animal do entrevistador) sabe que aqui em casa é o território dela, então ela fica aqui. Né?
\fala{E:} Já o gato macho, ele já sai mais.
\fala{A:} Não tem necessidade de andar. 
\fala{G:} É. Geralmente é os machos mesmo, né?
\fala{A:} É. É diferente. É.
\fala{A:} É difícil de ter território grande pra poder controlar. 
\fala{G:} Então, mas é isso. Eu acho que assim, tinha que conscientizar mais.
\fala{E:} É, até certo ponto.
\fala{G:} Pra não ficar multiplicando muito. Aí tem menos maltrato dos animais. Porque as pessoas às vezes colocam comida, mas outras vão lá e colocam veneno.
\fala{A:} Então...
\fala{E:} Uma das coisas que a gente pegou nas entrevistas, foi igual o caso do relato de uma criança que foi na ONG, que ela teve dois animais. Um foi envenenado e o outro, teoricamente, o pessoal disse que o telhado caiu em cima. Mas há um sério risco do pessoal ter jogado material de construção em cima do bicho pra matar.
\fala{G:} Sim.
\fala{E:} Né?
\fala{G:} É. Mas é por isso que eu falo. O Gabriel teve, ele teve dois cachorros.
\fala{G:} Um eles colocaram veneno pra ele. E o outro... Também parece que foi porque do jeito que ele morreu, ele adoeceu, muito de repente, e morreu também.
\fala{G:} Então assim, as pessoas elas... têm essas maldades. Então eu, particularmente, eu não gostaria de ter, de me apegar a alguma coisa que depois eu fosse sofrer e não tivesse condições de fazer nada.
\fala{G:} Ah, eu não gosto, gosto, acho bonito o cachorro e tudo, mas eu não pego por causa disso. Por não ter condições. Mas nem todo mundo.
\fala{G:} Igual você falou. Quando tá novo, tá beleza, não dá trabalho nenhum. Mas quando chega a fase adulta, a velhice, vai ter um trabalho igual um idoso mesmo.
\fala{G:} E que vai exigir cuidados. Então você tem que pensar hoje, não. Você tem que pensar amanhã também.
\fala{G:} O que vai fazer. Entendeu?
\fala{E:} Depois da tanta conversa(risos).
\fala{E:} A proposta partiu da existência e possibilidade do fato, ou problema de acordo como se interpreta a situação, das pessoas gostarem de ter animal de estimação.
\fala{E:} E geralmente também, quem tá na classe C, D e E. Não ter condições, justamente quando dá um problema. Ou ao cuidar do animal.
\fala{E:} A proposta que a gente teve foi o seguinte. Criar uma plataforma. Similar, basicamente, talvez.
\fala{E:} Igual comentei com a Aline(moradora entrevistada anteriormente). Criar uma espécie de Uber Pet de mão dupla. Ou seja, a pessoa obtem um serviço a ser disponibilizado.
\fala{E:} E pago por meio dos créditos da plataforma. Ou às vezes também pagando preço mais barato. Permitindo a possibilidade de conseguir cuidar do animal.
\fala{E:} E também se ela quiser prestar um serviço pra alguém, ela presta. Então, por exemplo: 
\fala{E:} "- Ah, eu não tenho condições de passear."
\fala{E:} "- E eu queria que alguém que estivesse próximo. Passeasse com o animal pra mim. Enquanto eu tô trabalhando."
\fala{E:} "- Ah, eu tô precisando levar o animal. Pra fazer uma consulta. Ou pra fazer algum procedimento veterinário."
\fala{E:} "- Mas eu não tenho condições de arcar com transporte."
\fala{E:} Ou às vezes também, por exemplo. A pessoa já tem a formação como veterinário.
\fala{E:} Mas aí ele não tem condições de abrir um escritório pra atender. Aí a pessoa basicamente iria até onde a pessoa tá precisando. Prestaria o serviço.
\fala{E:} E aí haveria a possibilidade da pessoa receber isso. Por meio dos créditos que vão ficar no programa. Ou também financeiramente.
\fala{E:} Esta é a nossa proposta. Aí eu vou abrir aqui pra mostrar. Aqui nesse caso agora.
\fala{E:} A gente tá mais na etapa de coletar feedback sobre a interface. Como é que tá o funcionamento dela. Se tá agradável.
\fala{E:} O que precisa melhorar. Ou se já tá no nível estável. Ou seja, coletar dados e informações.
\fala{E:} Aí o nome do programa é esse aqui, PetCuida. 
\fala{E:} Aqui no caso ele vai rodar uma simulação. Como seria um programa real.
\fala{E:} Na tela inicial ele tem aqui a opção da pessoa colocar a senha e o usuário. 
\fala{E:} E tem a opção de entrar pelo Google. 
\fala{E:} Tem essa simulação aqui pra parte de recuperar a senha.
\fala{E:} Que não tá montado, só sobe uma notificação. 
\fala{E:} Aqui a parte de criar a conta. Que aí segue o básico, colocar nome e senha, e confirmar.
\fala{E:} E aí faz o login. Tem a questão aqui de registrar. E de voltar ao menu.
\fala{E:} Nunca a gente conseguiu deixar pronto. Porque basicamente o programa foi feito pela graça de Deus em dois dias. E aí, ou seja, tá funcionando só a parte do registro Google.
\fala{E:} Aí nesse caso ele dá uma informação. E aí que a gente entra. Quando a pessoa entra.
\fala{E:} Como ele não guarda a informação. Então cada vez que ele faz um acesso ele volta. Então isso aqui é um processo padrão.
\fala{E:} Que vai acontecer quando a pessoa fizer o primeiro acesso. O tipo de animal disponível é geralmente cachorro, gato, pássaro. Ou então outro quando a pessoa tem um animal exótico.
\fala{E:} Aí no caso tem a opção de colocar o nome aqui. 
\fala{E:} Colocar aqui o velho e bom Rex da vida. Tem a opção de pegar uma foto do pet na galeria da pessoa.
\fala{E:} Ou também as vezes a pessoa se ele quiser pode bater a foto. 
\fala{E:} A gente só aperta ok.
\fala{E:} Isso aqui vai ficar armazenado como referência do animal. 
\fala{E:} No caso aqui a gente pode deixar por exemplo SRD. Que é sem raça definida.
\fala{E:} Tem aí na parte aqui de data de nascimento. O programa vai mostrar baseado em anos. Então por exemplo.
\fala{E:} Suponhamos que, nasceu em 2019. App vai pegar a data atual e reduzir pelo nascimento para achar a quantidade em anos. 
\fala{E:} No momento, como o app foi feito rápido. Não deu para poder deixar personalizado para o nosso jeito de ver calendário(estava em padrão americano). Por isso que ele está aqui em inglês um pouquinho.
\fala{E:} Então em geral, se a gente clicar aqui a pessoa escolhe o ano. Vou escolher no caso 2019. 
\fala{E:} A parte de mês ele se orienta para trás e para frente.
\fala{E:} Clica no dia. Aperta ok. Aí aqui ele já fez o cálculo aqui.
\fala{E:} Aqui a gente conclui. 
\fala{E:} E aqui esta é a interface principal. 
\fala{E:} Onde há as informações principais e serviços.
\fala{E:} A tela cumprimenta a pessoa. Mostra aqui a ideia geral das possibilidades que a pessoa tem. 
\fala{E:} A área que tem as informações do animal de estimação.
\fala{E:} A parte de processo de triagem. Para poder encaminhar a pessoa para o melhor atendimento possível. 
\fala{E:} A questão aqui para ver os serviços disponíveis ou quem pode atender.
\fala{E:} A parte do usuário mesmo. 
\fala{E:} Aqui a parte se a pessoa já quiser ver o que está acontecendo com o animal dele. 
\fala{E:} Aqui embaixo tem uma espécie de menu.
\fala{E:} Aqui embaixo a gente vê sobre o animal. 
\fala{E:} Aqui em cima tem o prontuário dele. O que ele já passou e o que está agendado.
\fala{E:} Aqui tem a opção de calendário de lembrete. Para a pessoa lembrar do que tem que fazer. Por exemplo.
\fala{E:} Tem que levar o animal para fazer uma consulta. Ou tomar vacina. Porque às vezes tem casos que a gente coletou a informação.
\fala{E:} Que a pessoa até esqueceu de vacinar o animal. 
\fala{E:} E tem a opção também no caso de compartilhar a informação do prontuário. Por meio de QR Code.
\fala{E:} Aqui é a fase de edição daqueles dados. 
\fala{E:} Então por exemplo. Ah não confundi era um gato.
\fala{E:} Pode também. Feito a alteração. E aperta aqui em salvar.
\fala{E:} Aqui embaixo é a parte de triagem. O que ele vai fazer perguntas para as pessoas. A ideia também.
\fala{E:} Depende do que for. Se precisar conversar com alguém. Aí fica mais fácil de orientar.
\fala{E:} A parte aqui de serviços disponíveis. E o que está em aberto. Quem está precisando.
\fala{E:} Tem aqui a ideia também de as pessoas serem moscadas no mapa. Mostrando quem está próximo ou não.
\fala{E:} Aí aqui do lado. Tem a parte que informa sobre o usuário. 
\fala{E:} O que ele fez.
\fala{E:} O que está para ser agendado. 
\fala{E:} A questão de suporte. Como eu falei com a senhora.
\fala{E:} Informação do saldo da pessoa. Dentro da plataforma. Os agendamentos.
\fala{E:} A parte de suporte técnico. 
\fala{E:} Aqui tem a opção de pegar a foto. Ou tirar a foto.
\fala{G:} Ou pegar da galeria. 
\fala{E:} No caso eu deixei assim. E ele já volta.
\fala{E:} Aqui do lado. Como todo programa exige. 
\fala{E:} A gente tem que deixar a informação de quem desenvolveu.
\fala{E:} Informar sobre os termos de uso. 
\fala{E:} E também sempre ter a informação da política de privacidade. 
\fala{E:} Por conta da lei geral de proteção de dados.
\fala{E:} Geralmente eles vão falar a tal da LGPD. 
\fala{E:} Com tudo assim que a senhora viu. O que a senhora achou?
\fala{G:} É um programa interessante. É interessante. E assim.
\fala{G:} É um programa bom. Só que as pessoas têm que conhecer esse programa também.
\fala{G:} Entendeu?
\fala{A:} Caso a pessoa se encontrar numa situação assim, de necessidade, ter algo assim.
\fala{G:} Pois, ajuda ter algo assim, ser uma coisa útil e oportuna. Algo que precisa ser passado para a frente.
\fala{A:} Pois muitas vezes, as pessoas tentam pesquisar na internet e achar.
\fala{E:} É, porque uma coisa igual a esta foi percebida. O Natanael(morador do bairro) comentou que, às vezes, pesquisa na internet, mas...
\fala{E:} Nem sempre dá certo, quando ocorre problema com os bichos. Ele tenta achar a solução. Pesquisa na internet.
\fala{E:} Mas às vezes a busca para o problema na internet pode levar a pessoa a errar a solução.
\fala{G:} Mas é interessante sim. Mas e se ele tivesse uma ajuda online?
\fala{E:} Mas a proposta dele é basicamente isso. Por isso que na parte de triagem. Ele tem a possibilidade da pessoa por exemplo.
\fala{E:} Clicar. E aí se tiver no caso, alguém.
\fala{E:} e esperamos conseguir progredir com o desenvolvimento.
\fala{E:} Ter ali um especialista para orientar o passo-a-passo do que se fazer. 
\fala{E:} Você pode ser encaminhado para nosso pessoal ou um usuário para ajudar.
\fala{E:} Para a pessoa não sair fazendo as coisas erradas. Ou seja, ter o tratamento correto. 
\fala{E:} E se resolver em tempo oportuno.
\fala{G:} E na verdade isso não gerar uma coisa assim que tipo assim. Aí eu vou ter porque tem quem cuida. Entendeu?
\fala{G:} Porque as pessoas elas precisam se conscientizar. Se tem elas têm que cuidar. 
\fala{G:} "Não é eu ter porque o Gabriel vai cuidar."
\fala{G:} "A Atanizia vai cuidar." 
\fala{A:} "Não, porque o animal é bonitinho."
\fala{G:} É. Porque é bonitinho. Entendeu?
\fala{G:} A pessoa tem que ter o animal, mas ter a consciência de que ela é que vai cuidar.
\fala{G:} Não dos outros.
\fala{G:} Que dá trabalho. 
\fala{A:} A maioria das pessoas que pega cachorro. Que pega gato.
\fala{A:} É porque acham bonitinho. 
\fala{G:} Sim. 
\fala{A:} Ou porque o filho quer. "Ah, porque é criança."
\fala{G:} Na hora que fica feio. 
\fala{A:} Na hora que adoece.
\fala{G:} Aí quer entregar para alguém. 
\fala{G:} É preciso assumir essa responsabilidade nesta hora também. Entendeu?
\fala{E:} Certo, desculpe interromper mas, em termos assim da interface. O que a sehora tem a dizer?
\fala{E:} Ficou prático para poder mexer? Se alguém ficou assim.
\fala{G:} Não. Ficou bom. Ficou fácil.
\fala{G:} Não teve. Não apareceu dificuldade. Está tranquilo.
\fala{G:} Não sei explicar, mas gostei.
\fala{A:} Ele está perguntando porque tem que coletar a opinião das pessoas.
\fala{E:} Coletar o feedback, ao mostrar a aplicação conforme informei antes da entrevista.
\fala{E:} Nós precisamos pegar entrevistas. Mostrar o programa. Explica um pouco do que foi a proposta.
\fala{E:} E verifica se, por exemplo. A interface precisa ser modificada. 
\fala{E:} Para ficar, digamos assim, mais intuitiva e atrativa. 
\fala{E:} Neste caso, o melhor possível para o cliente. 
\fala{E:} E também, ter um pouco de noção e de vivência de "Como é desenvolver um programa mesmo". \fala{E:} Porque ao criar um programa.
\fala{E:} O que a gente faz? 
\fala{E:} Primeiro você pesquisa a área que você quer atuar com aquele serviço.
\fala{E:} Depois é coletado a informação de: 
\fala{E:} - Pessoas
\fala{E:} - Sobre aquele problema. 
\fala{E:} - Para entender se é válido ou não o desenvolvimento.
\fala{E:} Depois começa-se o desenvolvimento. 
\fala{E:} E aí depois. É justamente iniciado o teste do programa.
\fala{E:} E coletando a informação do usuário. 
\fala{E:} Se está fácil de mexer. 
\fala{E:} O que precisa mudar.
\fala{E:} Apesar de que nessa conversa com a senhora, basicamente foram abordadas todas as etapas. 
\fala{E:} Desde a etapa inicial: Que era a parte de entender o problema. 
\fala{E:} Enquanto também coletava um pouco do feedback do usuário
\fala{E:} Antes de desenvolver até o ponto em que o programa já existe.
\fala{E:} Então se coleta mais feedback's como o da senhora.
\fala{E:} Acerca o programa que foi desenvolvido.
\fala{G:} Entendi.
\fala{E:} Então. Muito obrigado.
\fala{G:} De nada.
\end{turnos}
\section{Resultados da análise temática}
\subsection{Responsabilidade e compromisso do tutor}
A entrevistada associa a posse de um animal a uma responsabilidade contínua, comparável ao cuidado de um filho. A fala destaca que a decisão de ter um animal deve considerar despesas, doenças, envelhecimento e disponibilidade futura de cuidado. Esse tema funciona como limite normativo da própria plataforma: o serviço deve apoiar o tutor, mas não estimular a aquisição irresponsável.
\subsection{Barreiras econômicas e territoriais}
Os relatos apontam custos veterinários elevados, dificuldade de transporte adequado e concentração de serviços em Belo Horizonte. A distância, o preço de consultas e tratamentos e a ausência de atendimento local aparecem como barreiras de acesso. A plataforma é percebida como potencial mecanismo de mediação, especialmente quando conecta tutores, transporte, profissionais e serviços comunitários.
\subsection{Castração, abandono e maus-tratos}
A castração é apresentada como medida de prevenção à reprodução descontrolada e à sobrecarga de cuidadores. Também surgem relatos de abandono, envenenamento e possíveis maus-tratos. Como esses acontecimentos são relatos ou suspeitas dos participantes, este documento não os transforma em fatos verificados. Em pesquisa acadêmica, recomenda-se anonimização e cautela na divulgação.
\subsection{Utilidade percebida do PetCuida}
A entrevistada considera o programa interessante e afirma que ele deveria ser divulgado. As funcionalidades verbalizadas como relevantes incluem triagem, transporte, atendimento veterinário, lembretes, prontuário, rede de apoio e possibilidade de prestação de serviços mediante créditos ou remuneração.
\subsection{Facilidade de uso percebida}
O feedback final é positivo: a participante afirma que a interface ficou boa, fácil e sem dificuldade aparente. Essa evidência é limitada, pois a navegação foi conduzida pelo entrevistador e não houve execução independente de tarefas nem aplicação de escala padronizada, como a System Usability Scale (SUS).
\section{Limitações}
\begin{enumerate}[leftmargin=*]
\item A atribuição de Galdina às intervenções do terceiro falante é provisória.
\item A transcrição automática apresentou erros lexicais, nomes corrompidos e frases truncadas; os trechos incertos exigem nova escuta.
\item A amostra é pequena e não permite generalização estatística.
\item A relação familiar entre entrevistador e participante pode favorecer viés de cortesia.
\item A demonstração guiada não mede eficácia, tempo de tarefa, taxa de erro ou satisfação comparável entre versões.
\item Os valores, nomes de animais, locais e acontecimentos sensíveis devem ser confirmados antes de citação literal.
\end{enumerate}
\section{Recomendações para a próxima rodada}
\begin{enumerate}[label=R\arabic*.,leftmargin=*]
\item Confirmar previamente a identificação de cada participante e registrar essa informação no áudio.
\item Aplicar tarefas independentes: cadastrar um animal, consultar o prontuário, encontrar um serviço e solicitar transporte.
\item Perguntar separadamente sobre utilidade, confiança, clareza dos créditos, privacidade e intenção de uso.
\item Aplicar uma escala padronizada de satisfação ou usabilidade, além das perguntas abertas.
\item Criar explicação inicial sobre créditos, responsabilidades e limites da triagem, evitando interpretação do aplicativo como substituto de atendimento veterinário.
\item Validar juridicamente e eticamente o tratamento de dados pessoais e de informações de saúde animal.
\end{enumerate}
\section{Considerações finais}
A entrevista fornece evidência exploratória de que o PetCuida responde a problemas percebidos no cotidiano: custo, distância, transporte, castração e dificuldade de encontrar orientação. A recepção do protótipo foi favorável, sobretudo quanto à ideia de apoio e à facilidade percebida. O resultado, contudo, deve ser interpretado como feedback inicial, e não como validação estatística ou comprovação de usabilidade. A versão revisada da transcrição é adequada para análise qualitativa preliminar, desde que as marcações de incerteza sejam resolvidas por conferência final do áudio.
\appendix
\section{Nota de integridade da transcrição}
Os arquivos \texttt{.txt} e \texttt{.srt} entregues com este documento são versões semiliterais revisadas. Eles não substituem o áudio original. As alterações editoriais foram feitas para legibilidade e não devem ser apresentadas como transcrição palavra por palavra. Trechos entre colchetes indicam edição, incerteza ou necessidade de validação.
\end{document}
